\subsection{Proof of Theorem~\ref{thm:CDDP}}
We first give the intuition of the proof. We assume that each of the non-intersecting elements has high min-entropy. 
WLOG, consider an adversary corrupting $P_1$.
The view of the adversary will be 
  \[\view^{\Pbase}_{P_1}(A, B) := (A,c, h_1, \ldots, h_k). \]
%
As shown above, the sensitivity can be upper-bounded by a small value $s$. 

Unlike $\PC$, however, when we show the existence of sufficient noise from the remaining iterations, we need to take the additional leakage into consideration. 

First, since the hash functions are public, iterations are no longer independent of each other as needed by the analysis in Section~\ref{sec:curator}. We address this issue by employing the fact that each of the non-intersecting items has high min-entropy. In the random oracle model, as long as the adversary does not query hash function $h$ on some point $x$, $h(x)$ is uniformly random to the adversary.  Since the non-intersecting items have high min-entropy, the adversary is negligibly likely to query any of them to the hash functions, thus guaranteeing independence.  


\paragraph{Good iterations and Poisson Binomial distribution.}
Now, to see how the remaining iterations still hide the sensitivity even with the public hash functions, let  
$R = B \setminus I$. 
%
For the remaining $k-s$ iterations, the high min-entropy of each element in $R$ will jitter
the final count. In particular, consider the $j$th hash
function $h_j$ in the protocol (among the $k-s$ remaining iterations) and let
\[v_j^A = \min h_j(A),~~
v_j^I = \min h_j(I),~~
v_j^R = \min h_j(R).\]
%
Suppose $v_j^A = v_j^I$. Then, if $v_j^R \ge v_j^I$, the
min-hash $u_j^A$ of $A$ will be equal to the min-hash $u_j^B$ of $B$ (both of which are equal to $v_j^I$)
and the
final count $c$ will be incremented due to this $j$th iteration.
%
However, if $v_j^R < v_j^I$, then it will be $u_j^A \ne u_j^B$, and the
final count will not be incremented. This way, the distribution of $v_j^R$
will jitter the final count. 
%
The above discussion can be formalized into the following definition. 
\begin{definition}[$\theta$-good iteration]\label{def:theta-good}
Let $n_R = n_B - n_I$, we define 
$\good_\theta(h_j, A, I, n_B)$ to be true if and only if the following
holds: 
\[ \min h_j(A) = \min h_j(I), ~\mbox{and}~
    \min h_j(I) \in
    \left[1-\left(\frac{1}{2}+\theta\right)^{1/n_R}, 
        1-\left(\frac{1}{2}-\theta\right)^{1/n_R}\right].\]
\end{definition}
The second condition of the definition requires that $\min h_j(I)$ is somewhere in the middle (parameterized by $\theta \in \Theta(1)$) so that the distribution of $R$ (i.e., random $v_j^R$) may reduce the final count with a decent chance (and also keep the count with a decent chance).
%
As long as $n_I/n_A$ is a constant fraction, there are sufficiently many $\theta$-good iterations, although we lose some iterations.
In particular, if we let $k_g$ be the number of good iterations, we have $k_g = \Theta(k)$. 

With public hash functions and thereby $\min h_j(I)$ being leaked to the adversary, it turns out that the noise from the $k_g$ iterations follows a Poisson Binomial distribution, which is a generalization of a Binomial distribution where each trial has a different success probability. However, using the techniques of \cite{COK22}, we can still show that this distribution works as a good noise to hide the private data. 

%In particular, we  to upper-bound the privacy guarantee as those from a binomial mechanism (with rather worse parameters compared to $\PC$). 


 \subsubsection{Proof}
\label{apx:pubHashHighEntropy}
WLOG, we consider two neighboring inputs 
$(A, B) \mbox{ and } (A, \Bp).$
DDP for the case in which $\xs$ is added into $A$ can be shown symmetrically. We
prove the theorem by a hybrid argument. In particular, we define a slightly
different ideal functionality $\Pbase^{(1)}$ as follows:
%
\begin{itemize}
\item Let $\Pbase^{(1)}$ be the same as $\Pbase$ except that for each $x^B_i
\in B\setminus A$, each element in $\{y^B_{i,j}\}_j$ is chosen uniformly
at random from $[0,1]$.
\end{itemize}


We set up the following
hybrids.  We will argue that for any $\xs \in \U$ and over $(A, B, I)
\from \Delta_\PH$, it holds 
%
\begin{align*}
& (\view^{\Pbase}_{P_1}(A,B), I) \capprox
(\view^{\Pbase^{(1)}}_{P_1}(A,B), I) \\
& ~~~~~ \approx_{\ee, \dd} (\view^{\Pbase^{(1)}}_{P_1}(A, \Bp), I)
\capprox (\view^{\Pbase}_{P_1}(A, \Bp), I)
\end{align*}
%
for any constant $\ee > 0$ and for some $\dd \in \negl(\secparam)$, as
long as each element in $B \setminus I$ has high min-entropy.  

\iffalse
Thanks to the standard technique of the generic two-party secure
computation~\cite{JC:LinPin09}, the following holds: 
%
\[
\view^{\Pbase}_{P_1}(A,B), I) \capprox
(\view^{\Pbase}_{P_1}(A,B), I), ~~~
\view^{\Pbase}_{P_1}(A,\Bp), I) \capprox
(\view^{\Pbase}_{P_1}(A,\Bp), I). 
\]

Now we show $(\view^{\Pbase}_{P_1}(A,B), I) \capprox
(\view^{\Pbase^{(1)}}_{P_1}(A,B), I)$. 
%
\fi
Recall that the min-entropy of
each element $x^B_i$ with $i \in B\setminus A$ is at least $\kappa$.
Therefore, the probability that any adversary making at most
polynomially many oracle queries queries any $x^B_i$ is $\negl(\kappa)$.
Conditioned on the adversary not querying any such $x^B_i$, any
$y^B_{i,j}$ for $j \in [k]$ is chosen uniformly random from $\U$.
%
The same argument shows $\view^{\Pbase}_{P_1}(A,\Bp), I) \capprox
(\view^{\Pbase^{(1)}}_{P_1}(A,\Bp), I)$.
Therefore, we are left only to show
$(\view^{\Pbase^{(1)}}_{P_1}(A,B), I)  \aed  (\view^{\Pbase^{(1)}}_{P_1}(A,
\Bp), I)$.



\paragraph{DDP of $\Pbase^{(1)}$.}
We show $(A, I, h_1, \ldots, h_k, c) \aed (A,
I, h_1, \ldots, h_k, c_{+\xs})$, where $c$ is the final count from
$\Pbase^{(1)}(A, B)$ and $c_{+\xs}$ is the final count from $\Pbase^{(1)}(A, \Bp)$. 
%
We show how to leverage the uncertainties of $x^B_i
\in R = B \setminus A$ so that good iterations work like the needed noise to
guarantee DP.


\begin{lemma}\label{lem:good_event}
For any $A, I, n_B$ and $n_R = n_B - n_I$, we have
 \[p_\theta \stackrel{def}{=} \Pr_{h}[\good_\theta(h, A,
 I, n_B)] \ge
 \left(\left(\frac{1}{2}+\theta \right)^{\frac{n_A}{n_R}} -
 \left(\frac{1}{2}-\theta \right)^{\frac{n_A}{n_R}}\right) \cdot
 \frac{n_I}{n_A} ~~.\]
\end{lemma}
The proof is found in Section~\ref{apx:proof:good_event}.
This lemma shows that a random hash leads to a good iteration with
probability $p_\theta$, which is constant in our setting based on the assumption
about $n_A, n_I, n_R$.

Recall that $S_{\xs}$ was the random variable that represents the set of  
iterations $j$ such that the min-hash $u_j^B$ comes from $\xs$ when
$P_2$'s input is $\Bp$. From Lemma~\ref{lem:sen}, with overwhelming
probability $|S_{\xs}| = O(\lg\lg\kappa)$.  

Now, fix $A, I, \xs$ and $h_1, \ldots, h_k$ and let $G_{\theta}$ be the set of
iterations $j$ in which a $\theta$-good event takes place; i.e.,
%
$$G_{\theta} = \left\{ j \in [k]: \good_\theta(h_j, A, I, n_B) \right\}.$$

Let $K_{\theta} = G_{\theta} \setminus S_{\xs}$. 
The following lemma shows that the $\theta$-good events takes place
$\Theta(\secparam)$-many times, with overwhelming probability.  

\begin{lemma}\label{thm:good_hash}
Suppose $k = \Theta(\secparam)$, $n_B = \Omega(\secparam^2)$, and
$p_\theta \in \Theta(1)$. Let $s = |S_\xs|$. Then, 
we have 
\[ \Pr_{h_1, \ldots, h_k}\left[ |K_\theta| > \frac{2}{3} (k-s) p_\theta \right ] \ge 1 - \negl(\secparam).\] 
\end{lemma}
The proof is found in Section~\ref{apx:proof:good_hash}.



\paragraph{Our goal.}
For a set $W$, define $c_W \stackrel{def}{=} \sum_{j \in W} \Eq(u^A_j, u^B_j)$. Let $\overline K_\theta := [k] \setminus K_\theta$.
Note that the contributions to the final output can be divided into two parts:
\begin{itemize}
    \item  $c_{K_\theta}$:  The contribution from the iterations in $K_\theta$, which contains all the $\theta$-good iterations such that $\xs$ does not hash to the minimum across $\Bp$.
    
    \item $c_{\overline{K_\theta}}$:  The contribution from all the remaining iterations  
\end{itemize}
%

\iffalse
\paragraph{Fixing randomness.}
Now, fix $A, I, \xs$ and $h_1, \ldots, h_k$.
For sets $V, W$, let $Y_{V,W} \stackrel{def}{=}  \{ y^B_{i,j}: i \in V, j \in W\}$.
Recall that we are considering $\Pbase^{(1)}$ that satisfy the following condition: 
\begin{itemize}
\item For each $x^B_i \in B\setminus A$, each element in $\{y^B_{i,j}\}_j$ is chosen uniformly
at random from $[0,1]$.
\end{itemize}
Therefore, even after $(A, I, \xs, h_1, \ldots, h_k)$ is fixed, it holds that $Y_{R, [k]}$ are still uniformly distributed. 
Moreover, we additionally fix $Y_{R, \overline G_\theta}$, which implies that only $Y_{R, G_\theta}$ is uniformly random. 

Note that fixing 
$A, I, \xs, h_1, \ldots, h_k$ also allows us to determine which iterations are
$\theta$-good (i.e., $G_\theta$). 
%Thanks to Lemma~\ref{lem:sen},
%we can safely condition that the fixed hash functions are such that
%$|S_{\xs}|$ is smaller than or equal to $s =\lg\lg\kappa$.  

\paragraph{Small sensitivity.}
Given all this, we first show that there for $s = \Theta(\lg\lg\kappa)$, 
\[\Pr_{Y_{R, G_\theta}}\left[\left|c_{\overline K_\theta} -
c^{\xp}_{\overline
K_\theta}\right|> s \right] \le \negl(\kappa).\]
%
To show inequality, noting
that $\overline K_\theta = \overline G_\theta \cup S_{\xs}$, consider two cases: 

\begin{itemize}
\item For iteration $j\in \overline{G}_\theta \setminus S_{\xs}$, we have $u^B_j
= u^{\Bp}_j$. This is because $j \not \in S_{\xs}$ implies that $\xs$
doesn't lead to min-hash.  Therefore, such iterations don't
contribute to the difference of $c$ and $c^{\xp}$. 

\item For the rest iterations $j \in S_{\xs}$, we have $|S_{\xs}| \le s$ with overwhelming probability due to Lemma \ref{lem:sen}, which
shows the above equation. 
\end{itemize}
\fi


Essentially, for any final count $q$, we are interested in comparing the two
probabilities: 
%
\[\Pr[c_{\overline K_\theta} + c_{K_\theta} = q] \mbox{~and~}
\Pr[c^{\xp}_{\overline K_\theta} + c^{\xp}_{K_\theta} = q].\]

Following our discussion on sensitivity in Section \ref{sec:basic},  the
difference of $c_{\overline K_\theta}$ and $c^{\xp}_{\overline K_\theta}$
is upper-bounded by $s = O(\lg\lg \secparam)$. 
%
Note that we have $c_{K_\theta} = c^{\xp}_{K_\theta}$ because $j \in K_\theta$
implies $j \not \in S_{\xs}$. Therefore, we only need to analyze the single distribution of
$c_{K_\theta}$ as a noise and compare the following two probabilities:
%
\[\Pr[c_{K_\theta} = q] \mbox{~~and~}
\Pr[c_{K_\theta}+s = q].\]

\paragraph{Distribution of $c_{K_\theta}$.}
We have $c_{K_\theta} = \sum_{j \in K_\theta} c_j,$ where $c_j = {\sf Eq}(u_j^A, u_j^B)$. 
Note that since we have $j \in K_\theta$, a $\theta$-good event takes place in
iteration $j$, i.e., $\min h_j(A) = \min h_j(I)$. 

Let $\gamma_j = 1 - \min h_j(I)$. Note that the hash of each item $R$ is
randomly distributed in  $\Pbase^{(1)}$. Therefore, the probability that $c_j = 1$
is $(\gamma_j)^{n_R}$, in which case every hash of items in $R$ must be at least $\min h_j(I)$. 


%
Let $\eta_{-\theta} = 1/2-\theta$ and $\eta_{+\theta}= 1/2+\theta.$ Since $j$ is a good iteration, we have 
$(\gamma_j)^{n_R} \in [\eta_{-\theta}, \eta_{+\theta}].$
Therefore, letting $p_j = (\gamma_j)^{n_R}$, we have $c_j \sim \Ber(p_j)$, where
$\Ber$ denotes the Bernoulli distribution.  %
Since these Bernoulli distributions are independent from each other, 
can apply Lemma~\ref{lem:pb-dp} below to conclude that $C_{K_\theta} \aed
C_{K_\theta}+s$. 


\newcommand{\PB}{\mathsf{PB}}
For $j \in [n]$, consider  $c_j \sim \Ber(p_j)$.
% where
%$\Ber$ denotes the Bernoulli distribution.
%Let 
With $p_J = \{p_j\}_{j=1}^n,$
let $\PB(n, p_J)$ denote the distribution of $\sum_{j \in [n]} c_j$. This
distribution is called a Additive Poisson Binomial distribution. 

We conclude the proof by showing that the Additive Poisson Binomial distribution
with appropriate parameters satisfies the following DP-like property.  

\begin{lemma} \label{lem:pb-dp}
Consider an Additive Poisson Binomial distribution $\PB(n, p_J)$, where $n \in
\Omega(\kappa)$ and for each $p_j$, it holds that  $p_j \in [1/2 - \theta, 1/2 + \theta]$ where $\theta
\in (0, 1/2)$ is a constant independent of $\kappa$. Then, for any constant $\ee$ and $s  = \Theta(\lg\lg \kappa)$, 
there are $a, b \in [n]$ such that
\begin{itemize}
    \item For any $\ell \in [a,b], e^{-\ee} \le \frac{\Pr[\PB(n, p_J) =
    \ell]}{\Pr[\PB(n, p_J)+s = \ell]} \le e^{\ee}.$ 
    \item For any $\ell \not \in [a,b]$, $\Pr[\PB(n,p_J) = \ell] =
    \negl(\kappa)$ and $\Pr[\PB(n,p_J)+s = \ell] = \negl(\kappa)$.
\end{itemize}
\end{lemma}

The proof is found in Section~\ref{apx:proof:pb-dp}.

